\section{Polarization and the cost of a complementary modulus}\label{sec:tariff}

Fix a real parent and an original arithmetic itinerary. We count integers
$u$ in a bin $U/2\le u\le U$, not factorizations chosen before seeing the
parent. Write $G_u=(1/u)\mathbb Z/\mathbb Z$. Suppose that $B_u\subset
G_u$ is nonempty, exactly invariant under multiplication by $3$, and
$B_u+B_u$ is covered by at most $U^\theta$ translates of $B_u$.
The argument also applies at the terminal node, with $u$ the residual
modulus. All entropies and logarithms in this section are in nats.

\begin{lemma}[Polarization in prime digits]\label{lem:digits}
Fix $0<c,\delta,\rho<1$. There exist integers $m,R$, depending only on
these parameters, with the following property. Let $3\nmid u$, let
$\varnothing\ne B\subset G_u$ be invariant under multiplication by $3$,
and suppose $B+B$ is covered by $M$ translates of $B$. Use the Chinese
remainder isomorphism from $G_u\simeq\mathbb Z/u\mathbb Z$ to its
prime-power coordinates. In any fixed order of their base-prime digits
that places lower digits before higher digits in each coordinate, let
$h_i$ be the entropy of digit $i$ of a uniform $X\in B$, conditional on
all preceding digits. Denote its base by $r_i$. For any selection $I$
of digits satisfying
\[
 r_i>R,\qquad \operatorname{ord}_{r_i}(3)\ge r_i^c,
 \qquad h_i\ge\delta\log r_i ,
\]
one has
\[
 \sum_{i\in I}(\log r_i-h_i)
 \le \rho\sum_{i\in I}\log r_i+(m-1)\log M .
\]
\end{lemma}

\begin{proof}
Take $m$ independent copies $X_1,\ldots,X_m$ of $X$, and, independently,
choose $t_1,\ldots,t_m$ uniformly and independently modulo
$\operatorname{ord}_u(3)$, using least nonnegative representatives.
The case $u=1$ is empty. Put $Y=\sum_l3^{t_l}X_l$, and let
$e_i=H(Y_i\mid Y_{<i},\boldsymbol t)$. In bounding $e_i$ below we may
instead condition on $\boldsymbol t$ and the $m$ individual input
histories $X_{l,<i}$: these determine every preceding output digit.
Indeed, for a digit at level $j$ (the units digit is level $0$) of
an $r$-primary coordinate, let $b_l$ be its known residue modulo $r^j$
in $[0,r^j)$, and $v_l$ its current digit. The output digit is
\[
 \left\lfloor r^{-j}\sum_l3^{t_l}b_l\right\rfloor+
 \sum_l(3^{t_l}\bmod r)v_l \pmod r .
\]
The first term is known. The remaining terms are independent conditional
on the individual histories and the times. Adding independent variables
on a finite group cannot lower the entropy of any one summand
(condition on the other summands). Thus $e_i\ge h_i$ for every $i$.

We claim also $e_i\ge(1-\rho)\log r$ whenever the hypotheses for $I$
hold. Write $\mu_z$ for the conditional digit law given a single history
$z$, and $a_z$ for its mass on atoms of size at most $r^{-\delta/4}$.
Separating these atoms from their complement gives
\[
 H(\mu_z)\le\log 2+(\delta/4)\log r+a_z\log r
\]
since the complement has at most $r^{\delta/4}$ elements.
Hence $\mathbb E a_z\ge\delta/2$ for large $r$. Flag a copy if its
current atom is in that small-atom set and $a_z\ge r^{-\delta/8}$.
The probability of a flag is at least $\delta/3$, and conditional on
its history and positive flag a copy has maximum atom at most
$r^{-\delta/8}$. The flags of the copies are independent and do not
depend on the times. We may lower bound the output entropy by
additionally conditioning on all these flags.

We use the prime-field subgroup exponential-sum estimate of
Bourgain, Glibichuk and Konyagin \cite{BGK}: for some $b=b(c)>0$ and
$C=C(c)$, if $\Gamma\subset\mathbb F_r^\times$ is a subgroup of size
at least $r^c$, the absolute value of the average of any nontrivial
additive character over $\Gamma$ is at most $C r^{-b}$.
For a flagged conditional law $\nu$ and $v\ne0$, expand the squared
Fourier coefficient with two independent samples from $\nu$.
The diagonal has mass at most $r^{-\delta/8}$; off the diagonal
$3^{t_l}\bmod r$ is uniform on $\langle3\rangle$, so
\[
 \mathbb E_{t_l}|\widehat\nu(3^{t_l}v)|^2
 \le r^{-\delta/8}+Cr^{-b}\le r^{-b'}
\]
for some $b'>0$ and large $r$. Fix $J$ with $Jb'\ge2$.
Conditional on histories and flags with at least $J$ positive flags,
independence of the times and Parseval give mean conditional collision
probability for the output digit at most
$r^{-1}(1+(r-1)r^{-Jb'})$. The known shift can depend on all the times
without changing this probability; the remaining, possibly unflagged,
Fourier factors have modulus at most one. Entropy is at least minus the
logarithm of collision probability, so Jensen gives mean conditional
entropy at least $\log r-\log(1+r^{-1})$ in this case. Choose $m$ so that
the probability of fewer than $J$ successes in $m$ Bernoulli trials of
parameter $\delta/3$ is at most $\rho/2$, then take $R$ sufficiently
large. This proves the claim. Only independence \emph{between copies}
has been used; the clocks of different primes in one copy may coincide.

Finally $Y$ belongs to the $m$-fold sumset of $B$. Iterating the
translate cover shows that this sumset has at most $M^{m-1}|B|$
elements. Consequently
\[
 0\le\sum_i(e_i-h_i)=H(Y\mid\boldsymbol t)-\log|B|
 \le (m-1)\log M .
\]
Applying the two lower bounds for $e_i$ proves the lemma.
\end{proof}

\begin{lemma}[Prime-base catalogue]\label{lem:basecatalog}
Fix $P_0$, $\xi>0$, and $1/2<s<1$. For $d$ write $d_{\rm big}$ for the
factor composed of prime powers at primes exceeding $P_0$. Consider
either all primes $r\ne3$ for which $d=\operatorname{ord}_r(3)$
satisfies $d_{\rm big}=1$, or, for one fixed $L>1$ having only primes exceeding
$P_0$, all those for which $d_{\rm big}\mid L$ and $d_{\rm big}\ge L^\xi$.
In either catalogue the number of squarefree products of the allowed
bases not exceeding $V\ge1$ is $O_{s,\xi,P_0}(V^s)$, uniformly in $L$.
\end{lemma}
\begin{proof}
At a fixed order $d$ there are $O(d)$ possible primes since their
product divides $3^d-1$. Each is $1\bmod d$. The $j$th smallest is
at least $1+jd$, whence
\[
 \sum_{\operatorname{ord}_r(3)=d}r^{-s}
 \ll_s d^{-s}\sum_{j\ll d}j^{-s}\ll_s d^{1-2s}.
\]
Write $a=2s-1>0$. The sum of $w^{-a}$ over integers with all prime
factors $\le P_0$ is $\prod_{p\le P_0}(1-p^{-a})^{-1}$, with the
product taken over primes. In the second case this constant is
multiplied by at most $\tau(L)L^{-a\xi}=O_{a,\xi}(1)$, by the divisor
bound. Thus $\sum_r r^{-s}$ is uniformly bounded and so is
$\prod_r(1+r^{-s})$. Multiplying the latter by $V^s$ gives the count.
\end{proof}

\begin{proposition}[Adaptive complementary count]\label{prop:tariff}
Fix a large-drop node or a terminal node, its original arithmetic
labels (in particular $n=\operatorname{ord}_q(3)=T^{o(1)}$), and
its actual preceding observations. Take $T^\eta\le U\ll T$ with
fixed $\eta>0$. For every $\sigma>0$, with the covering exponent
$\theta$ sufficiently small and the tolerance in Lemma~\ref{lem:lcm}
sufficiently small, the number of eligible moduli $u\in[U/2,U]$ with
$|B_u|\ge U^\beta$ satisfies, uniformly for $0\le\beta\le1$,
\[
 \#\{u\}\ \ll\ T^{\xi+o(1)}U^{(1-\beta)/2+\sigma}.
\]
At a terminal node the factor $T^\xi$ is omitted. Here $P_0$ may be
chosen after the covering and quotient parameters; the implied
constant may depend on its final value.
\end{proposition}

\begin{proof}
Recall the arithmetic information used at a single node. At a prime
drop the part of the entering lag supported on primes $>P_0$ is $L$.
Outside primes dividing $n$, the low-order part lost has product at
most $T^\xi$ by the itinerary choice; all other lost bases are in
the second catalogue of Lemma~\ref{lem:basecatalog}. At the terminal
all orders divide the terminal lag, so the first catalogue applies.
These are statements about the original itinerary. Indeed, if
$r^e\mid q$ and $r\nmid 3n$, the order at $r^e$ equals the order at
$r$: their quotient is an $r$-power and divides $n$. Such a component
leaves the lag gcd in its entirety or not at all. Partial losses over
$r\mid n$ need not satisfy this assertion.

Choose a large integer $A$, small $c,\delta,\rho>0$ and
$1/2<s<1$ close to $1/2$. Let $m,R$ be as in Lemma~\ref{lem:digits}.
Put aside, in the following order, the parts of $u$ listed here.
Their products, taken separately, have the indicated number of
possible integer names; their heights need not themselves be small.
\begin{enumerate}
\item The part on primes dividing $n$ or at most $R$: $T^{o(1)}$ names.
For each fixed $v>0$ the number is at most
$U^v\prod_{r\mid n\ \mathrm{or}\ r\le R}(1-r^{-v})^{-1}$.
The logarithm of the product is $O_{v,R}(1+\log(2n))=o(\log T)$.
Taking arbitrarily small fixed $v$ gives the assertion.
\item The remaining low-order loss of the itinerary: at most $T^\xi$
names by its height, and no such part is needed at the terminal.
\item All remaining full powers $r^e$ with $e\ge A$: $O_A(U^{2/A})$
names, since
$\prod_r(1+\sum_{e\ge A}r^{-2e/A})<\infty$.
\item All remaining bases with $\operatorname{ord}_r(3)<r^c$,
including their powers: $O_c(U^{3c})$ names. Indeed, at most
$O(Y^{2c})$ such bases lie below $Y$, by summing $O(d)$ over
$d\le Y^c$. Dyadic summation gives $\sum_r r^{-3c}<\infty$
over these bases, so the full geometric Euler product converges.
\end{enumerate}
For the rest, call a full block $r^e$ active if \emph{every} digit
has $h_i\ge\delta\log r$ for $X$ uniform on $B_u$, using a fixed
order for that $u$. Put the product of these whole blocks in $D$, and
put $K=u/D$. Empty products are allowed. The complementary factor
$K$ includes all parts put aside. Lemma~\ref{lem:digits} and
$M\le U^\theta$ imply that the sum of the entropy deficits of the
active digits is at most
$\alpha\log U$, where $\alpha=\rho+(m-1)\theta$.
Conditioning only on previous active digits instead of on the full
histories can only increase the corresponding entropies. The
projection of $B_u$ on its active blocks therefore has at least
$D U^{-\alpha}$ elements. Since $(D,K)=1$, multiplication by $K$
kills exactly the complement and is invertible on the active blocks.
In particular
\[
                 |K\cdot B_u|\ge D U^{-\alpha}.
\]
For each \emph{fixed} integer $K$, use Lemma~\ref{lem:lcm} with
tolerance $\lambda>0$: its sum of $|K\cdot B_{KD}|$ over eligible
$D$ is $\ll (U/K)U^\lambda T^{o(1)}$. Since $D\ge U/(2K)$, there
are at most $O(U^{\alpha+\lambda}T^{o(1)})$ choices of $D$ in
our selection, including when $K=1$ or $D=1$.

It remains to count $K$, which depends on the set $B_u$. Every
remaining inactive block has a digit with deficit at least
$(1-\delta)\log r$. Since all deficits are nonnegative and sum
to $\log u-\log|B_u|$, the product of these distinct bases satisfies
\[
    \sum_{r\ {\rm remaining\ inactive}}\log r
           \le \frac{(1-\beta)\log U}{1-\delta}.
\]
For $1\le e<A$ put $V_e$ equal to the least power of two not below
the product of these bases of exponent $e$ (use $1$ for an empty
product). There are $O_A((1+\log U)^A)$ vectors of bins. In any
nonempty vector,
$\sum_e\log V_e\le A\log2+(1-\beta)\log U/(1-\delta)$.
For each $e$ the eligible base products are squarefree subproducts
of the \emph{original} node catalogue, so at most
$O_{s,\xi,P_0}(V_e^s)$ can occur. Raising a chosen product to the
known exponent $e$ recovers that part of $K$. This counting argument
uses neither a newly valid itinerary after deletion nor independence
among the factors; the bins are a device for a union bound at this
parent, not extra labels on the stationary point distribution.

Multiplication of the counts for the parts put aside, the remaining
inactive names, and the choices at fixed $K$ gives the asserted
estimate with exponent on $U$ at most
\[
 \frac{s(1-\beta)}{1-\delta}
       +\frac2A+3c+\rho+(m-1)\theta+\lambda ,
\]
besides $T^{o(1)}$ and the possible $T^\xi$. First choose
$A,c,\delta,\rho,s$ to make the excess over $(1-\beta)/2$
small uniformly in $\beta$; next fix $m,R$, then $\theta$ and
$\lambda$. The large time-prime cutoff needed to realize the
cover and Lemma~\ref{lem:lcm} can be taken afterwards.
\end{proof}
